\documentclass[a4paper]{article}

\usepackage{cmap}
\usepackage[T2A]{fontenc}
\usepackage[utf8]{inputenc}
\usepackage[russian]{babel}
\usepackage{indentfirst}
\usepackage[hidelinks]{hyperref}
\usepackage{pgfplots}
\usepackage{amsmath}
\usepackage{amssymb}
\usepackage{circuitikz}

\title{Отчёт по лабораторной работе\\<<Опыт Франка-Герца>>}
\author{Оськин И.Д. \and Можаров А.Р.}

\begin{document}

\maketitle

\section*{\centering Теоретическая часть}

\subsection*{Постулаты Бора}

В 1911 году на основании экспериментов Резерфордом была предложена планетарная модель атома,
однако классические механика и электродинамика не могли объяснить устойчивость атома
и его спектр. Для этого в 1913 году Нильсом Бором были предложены постулаты
квантовой теории строения атома, которые в последствии были названы {\it постулатами Бора}
\begin{enumerate}
    \item Атомы могут пребывать длительное время только в состояниях с определёнными энергиями,
    образующими дискретный ряд $E_0$, $E_1$, \ldots При нахождении атома в этих состояниях, атом не поглощает
    и не излучает, а излучение энергии атомом происходит только скачкообразно и строго между
    указанными энергетическими уровнями.
    \item При переходе из одного состояния с энергией $E_m$ в другое с энергией $E_n$
    может поглотиться или излучиться строго определённая порция электромагнитной энергии
    (называемая {\it квантом} или {\it фотоном}). Энергия одного кванта излучения
    в зависимости от частоты излучения может быть выражена как
    $$
        h\nu = \hbar\omega = E_m - E_n
    $$
    где $h$, $\hbar$ --- постоянные планка.
\end{enumerate}
Непосредственное подтверждение постулатов Бора было получено в том же 1913 году
в результате экспериментов, выполненных Д. Франком и Г. Герцем.

\subsection*{Установка в опыте Франка-Герца}

В опыте Франка-Герца используются столкновения электронов с атомами газа.
Схему опытной установки см. на рис. \ref{Установка}.
Баллон триода (газовой лампы с тремя электродами: анодом, катодом и сеткой)
заполнен исследуемым газом. Катод является источником электронов,
а анод, соответственно приёмником. Третий электрод --- сетка --- играет следующую роль:
между сеткой и катодом создаётся положительное напряжение
$$
    U_\text{уск.} = \varphi_\text{сет.} - \varphi_\text{кат.}
$$
называемое {\it ускоряющим потенциалом} (на рис. \ref{Установка} измеряется вольтметром $U_\text{у.}$),
а между анодом и сеткой, как правило противоположного знака, напряжение
$$
    U_\text{зад.} = \varphi_\text{ан.} - \varphi_\text{сет.}
$$
называемое {\it потенциалом задержки} (на рис. \ref{Установка} измеряется вольтметром $U_\text{з.}$).

\begin{figure}
    \centering
    \begin{circuitikz}[european resistors, american inductors]
        \draw (0, 0) to [rmeter, t=$I_\text{а.}$, -*] (2, 0) to [rmeter, t=$U_\text{з.}$, -*] (2, -2) -- (4, -2) to [pR, wiper pos=0.5, *-, name=R1] (4, 0) -- (5, 0) to [short, -o] (5, -0.75);
        \draw (2, 0) -- (3, 0) -- (3, -1) -- (R1.wiper);
        \draw (5, -1.25) to [short, o-] (5, -2) -- (4, -2);
        \draw (2, -2) -- (2, -2.25) to [rmeter, t=$U_\text{у.}$, -*] (2, -4.25) -- (4, -4.25) to [pR, wiper pos=0.5, *-, name=R2] (4, -2.25) -- (5, -2.25) to [short, -o] (5, -3);
        \draw (5, -3.5) to [short, o-] (5, -4.25) -- (4, -4.25);
        \draw (3, -2) to [short, *-] (3, -3.25) -- (R2.wiper);
        \draw (0, -2) node[triode, fullcathode, name=T, xscale=-1]{};
        \draw (0, 0) -- (T.anode);
        \draw (2, -2) -- (T.control);
        \draw (T.cathode) to [short, -*] (0.28, -3.25) -- (1, -3.25) -- (1, -4.25) -- (2, -4.25);
        \draw (-0.28, -2.95) -- (-0.28, -3.25) to [short, -o] (-0.28, -4.25);
        \draw (0.28, -3.25) to [short, -o] (0.28, -4.25);
    \end{circuitikz}
    \caption{Принципиальная схема экспериментальной установки}
    \label{Установка}
\end{figure}

\subsection*{Качественный анализ эксперимента}

В ходе выполнения эксперимента снимается анодно-сеточная характеристика
(т.е. зависимость анодного тока $I_\text{а.}$ от ускоряющего потенциала $U_\text{у.}$
при постоянном задерживающем потенциале $U_\text{з.}$),
при этом будет наблюдаться следующее (будем увеличивать $U_\text{у.}$)
\begin{enumerate}
    \item При $U_\text{у.} < U_\text{з.}$ электроны не успевают приобрести достаточную энергию для возбуждения атома,
    а также не обладают достаточной энергией, чтобы пройти всё расстояние от сетки до анода
    (однако некоторые случайные электроны такой энергией обладать могут и по мере увеличения
    $U_\text{у.}$ их количество будет расти, поэтому будет наблюдаться лёгкий рост анодного тока).
    \item При $U_\text{у.} > U_\text{з.}$ электроны будут иметь энергию,
    достаточную для преодоления задерживающего потенциала, поэтому анодный ток будет расти.
    \item При достижении {\it резонансного/критического потенциала}
    $$
        U_1 = \dfrac{E_1 - E_0}{e}
    $$
    электроны уже будут обладать энергией для возбуждения атомов, поэтому анодный ток $I_\text{а.}$
    начнёт спадать, т.к. электроны, потерявшие энергию при столкновениях с атомами,
    уже не смогут преодолеть задерживающий потенциал.

    Атомы в возбуждённом состоянии $E_1$ буду стремиться вернуться в основное состояние
    и при этом процессе будут испускать фотон с частотой
    $$
        \nu = \dfrac{E_1 - E_0}{h}
    $$

    При дальнейшем увеличении ускоряющего потенциала,
    область накопления электронов будет смещаться ближе к аноду.
    \item При $U_\text{у.} > U_1 + |U_\text{з.}|$ электроны,
    испытавшие соударение, начнут приобретать энергию, достаточную как для соударения,
    так и для преодоления задерживающего потенциала.
    \item Начиная с $U_\text{у.} > U_2$ (где $U_2 = 2U_1$) процессы будут аналогичны тем,
    что описаны в предыдущих двух пунктах (и т.д. $U_n = nU_1$).
\end{enumerate}
Таким образом опыт Франка-Герца подтверждает оба постулата Бора.

\subsection*{Ионизационный потенциал}

Рассмотрим ситуацию, когда $|U_\text{з.}| > U_\text{у.}$. В этом режиме электроны
не достигают анода, и анодный ток может быть вызван только
положительными носителями заряда (ионами газа).

Когда $U_\text{у.}$ достигает значения $E_\text{ион.}$ происходит ионизация
$$
    A + e^- + E_\text{ион.} \rightarrow A^+ + 2e^-
$$
и на одно-сеточной характеристике наблюдается скачок анодного тока в
противоположном обычному направлении.

\section*{\centering Практическая часть}

\subsection*{Определение резонансного уровня}

Измерения, проведённые в ходе определения резонансного уровня при $U_\text{з.} = 7,5$ В и $U_\text{нак.} = 3,5$ В,
приведены на рис. \ref{Определение резонансного уровня}.

Определение резонансного потенциала производится из его физических свойств.
Точки $U_n$ будут являться максимумами анодно-сеточной характеристики. Тогда
$$
    U_1 = \dfrac{U_n - U_1}{n - 1} \approx 41,5 \text{ В } - 21 \text{ В } = (20,5 \pm 0,25) \text{ В}
$$
Разность уровней $E_1 - E_0$ может быть определена из соотношения
$$
    E_1 - E_0 = e U_1 = (20,5 \pm 0,5) \text{ эВ } \approx (3,28 \pm 0,04) \cdot 10^{-18} \text{ Дж}
$$

\begin{figure}[h]
    \centering
    \begin{tikzpicture}
        \begin{axis}[
            title={Анодно-сеточная характеристика},
            xlabel={Ускоряющий потенциал $U_\text{у.}$, [0-75] В},
            ylabel={Анодный ток $I_\text{а.}$, [0-100] мкА},
            grid=major,
            xmin=-5, xmax=155,
            ymin=-0.1, ymax=1.1
        ]
        
        \addplot[mark=*, line width=1pt, blue] table[x index=0, y index=1]{Graph-1.txt};

        \end{axis}
    \end{tikzpicture}
    \caption{Определение резонансного уровня}
    \label{Определение резонансного уровня}
\end{figure}

\subsection*{Определение ионизационного потенциала}

Измерения, проведённые в ходе определения ионизационного потенциала,
приведены на рис. \ref{Определение ионизационного потенциала}.

Потенциал ионизации определяется как потенциал, в котором происходит
резкий скачок анодно-сеточной характеристики, тогда
$$
    U_\text{ион.} \approx (24,5 \pm 0,25) \text{ В}
$$
Откуда энергия ионизации
$$
    E_\text{ион.} = e U_\text{ион.} \approx (3,92 \pm 0,04) \cdot 10^{-18} \text{ Дж}
$$
Также может быть посчитано изменение энергии электрона на длине свободного пробега
$$
    \Delta W_\lambda \approx e \cdot \lambda \cdot \dfrac{dU}{dn} \approx e (U_\text{ион.} - U_1) = (4 \pm 0,5) \text{ эВ } \approx (6,4 \pm 0,8) \cdot 10^{-19} \text{ В}
$$

\begin{figure}[h]
    \centering
    \begin{tikzpicture}
        \begin{axis}[
            title={Анодно-сеточная характеристика},
            xlabel={Ускоряющий потенциал $U_\text{у.}$, [0-75] В},
            ylabel={Анодный ток $I_\text{а.}$, [0-1] мкА},
            grid=major,
            xmin=32, xmax=55,
            ymin=-0.2, ymax=1.2,
            legend pos=north west
        ]
        
        \addplot[mark=*, line width=1pt, blue] table[x index=0, y index=1]{Graph-2-1.txt};
        \addlegendentry{$U_\text{з.} = 35 \text{В}$}
        \addplot[mark=*, line width=1pt, green] table[x index=0, y index=1]{Graph-2-2.txt};
        \addlegendentry{$U_\text{з.} = 40 \text{В}$}
        \addplot[mark=*, line width=1pt, red] table[x index=0, y index=1]{Graph-2-3.txt};
        \addlegendentry{$U_\text{з.} = 45 \text{В}$}

        \end{axis}
    \end{tikzpicture}
    \caption{Определение ионизационного потенциала}
    \label{Определение ионизационного потенциала}
\end{figure}

\section*{\centering Дополнительные вопросы}

\begin{enumerate}
    \item {\it Как изменится вид анодно-сеточной характеристики в условиях опыта,
    если рабочий газ с заметной вероятностью захватывает электроны?
    Можно ли в этом случае определить разности энергетических уровней атомов и их ионизационные потенциалы?}

    Если рабочий газ с заметной вероятностью захватывает электроны
    (образует отрицательные ионы),
    то это приведет к уменьшению анодного тока $i_a$,
    поскольку меньше электронов сможет достичь анода1.

    \item {\it Будет ли в условиях опыта соблюдаться равенство $U_1 = U_2 - U_1$?
    Потенциал $U_1$ или разность $U_2 - U_1$ следует считать резонансным потенциалом и почему?}

    Первый критический потенциал $U_1$ соответствует энергии $eU_1 = E_1 - E_0$,
    необходимой для первого возбуждения атома.
    При потенциале $U_\text{у.} \geqslant 2U_1$ электроны могут испытать два неупругих столкновения,
    что приводит к появлению второго провала на характеристике при $U_\text{у.} \approx U_2$.
    Если предположить, что энергия, потерянная при каждом неупругом столкновении,
    одинакова и равна энергии первого возбуждения, то $U_2 \approx 2U_1$.
    Следовательно, $U_2 - U_1 \approx 2U_1 - U_1 = U_1$.

    \item {\it Известно, что в результате термоэлектронной эмиссии электроны вылетают из катода,
    имея различные скорости. Повлияет ли это обстоятельство на точность определения резонансного потенциала?}

    Да, то, что электроны вылетают из катода, имея различные начальные скорости
    (и, следовательно, энергии), повлияет на точность определения резонансного потенциала.
    Разброс начальных скоростей приведет к разбросу электронов по энергиям при достижении ими области неупругих соударений.
    Этот разброс будет уменьшать крутизну падающего участка характеристики и расширять максимумы.
    Расширение максимумов и уменьшение крутизны затрудняют точное определение точек $U_1$ и $U_\text{мин.}$,
    что, в конечном итоге, снижает точность определения резонансного потенциала.

    \item {\it Как изменятся анодно-сеточные характеристики с увеличением густоты сетки лампы?}
    
    Тогда $U_\text{у.}$ находится в области первого минимума ($U_\text{у.} \approx U_\text{мин.}$),
    электроны, потерявшие энергию при неупругих столкновениях, летят на анод,
    минуя область неупругих соударений.
    Если воронки неглубокие (что может произойти при увеличении густоты сетки),
    меньше электронов будет терять энергию вблизи "дна" воронки.
    Если же густота сетки станет очень большой, приближаясь к сплошной,
    это может привести к тому, что электроны при $U_\text{у.} \approx U_\text{мин.}$
    смогут попасть на анод, и ток в минимуме может доходить до нуля.

    \item {\it Почему выполнение условия $|U_\text{з.}| > U_\text{у.}$
    необходимо для наблюдения скачка анодного тока при ионизации?}

    Это условие необходимо, чтобы электроны, испущенные катодом,
    не достигали анода до тех пор, пока не произойдет ионизация.
    
    \item {\it С чем связано изменение знака анодного тока при наблюдении ионизации?}
    
    Это связано с изменением знака частиц (электроны $\rightarrow$ ионы), которые обеспечивают этот ток.
\end{enumerate}

\end{document}